% ============================================================================
% S\'eCURIT\'e GAME-BASED ET R\'eDUCTIONS FORMELLES
% source : annexes/Gamebased.tex
% ============================================================================

\chapter{S\'ecurit\'e game-based et r\'eductions formelles pour la famille SOTP}
\label{chap:annexe_gamebased}

\begin{sloppypar}
Le cadre CLO d\'efinit $\Conf$, $\Rel$ et $\Opp$ avec des termes issus de la théorie de l'information, \`a savoir les entropies conditionnelles, l'information
mutuelle et la distance de variation totale~\cite{MINKA2026CLO}. Cette approche
permet des r\'esultats g\'en\'eraux (th\'eor\`eme du goulot d'\'etranglement
contextuel, Trilemme CRO conditionnel), mais ne pr\'ecise pas directement
\emph{ce qu'un adversaire calculatoirement born\'e peut ou ne peut pas
faire}. En cryptographie computationnelle, la s\'ecurit\'e s'exprime
classiquement via des \emph{jeux} (games) entre un challenger et un
adversaire probabiliste polynomial-time (PPT).

\begin{cmtEtape}{Pourquoi cette annexe.}
Le cadre CLO raisonne en bits et en entropies, alors que la cryptographie classique
raisonne en jeux et en avantages d'adversaires. Cette annexe construit
le pont entre les deux en capturant chaque m\'etrique du cadre CLO
par un jeu cryptographique, puis en reliant chaque jeu \`a une hypoth\`ese
de duret\'e standard (PRG, DLOG, JL, Groth16).
\end{cmtEtape}

Cette annexe \'etablit le pont formel entre les deux mondes en trois \'etapes.
\begin{enumerate}
\item D\'efinir des jeux cryptographiques (\emph{IND-CONF},
      \emph{BIND-REL}, \emph{EXT-OPP}) capturant op\'erationnellement
      $\Conf$ (dans sa version calculatoire), $\Relver$ et $\Opp$
\item Prouver des \emph{r\'eductions} formelles \'etablissant que, si un adversaire PPT
      gagne l'un de ces jeux, il r\'esout un probl\`eme suppos\'e difficile
      (PRG, logarithme discret, JL, Groth16)
\item D\'eriver une minoration game-based du produit CLO
      (Corollaire~\ref{cor:ann_cro_from_games}), \`a l'\'echelle op\'erationnelle
      (confidentialit\'e calculatoire, fiabilit\'e de v\'erification), qui ne se
      compare pas directement au Trilemme
      (Remarque~\ref{rem:ann_pas_de_contradiction})
\end{enumerate}

\begin{cmtEtape}{Trois jeux, trois m\'etriques.}
\textbf{IND-CONF} capte la confidentialit\'e en exigeant que l'adversaire ne puisse
distinguer les preuves de deux messages. \textbf{BIND-REL} capte la
fiabilit\'e, puisque l'adversaire ne doit pas pouvoir ouvrir une preuve sur
deux messages distincts. Enfin, \textbf{EXT-OPP} capte l'opposabilit\'e,
l'interpr\`ete institutionnel devant extraire le bon verdict depuis la
preuve.
\end{cmtEtape}

\noindent\textbf{Convention.} Tout adversaire $\mathcal{A}$ est
probabiliste polynomial-time (PPT) en $\lambda$ sauf mention contraire.
Le challenger est not\'e $\mathcal{C}$. Une r\'eduction construit un
algorithme $\mathcal{B}$ (r\'educteur) qui exploite $\mathcal{A}$ pour
r\'esoudre un probl\`eme difficile.
\cmt{Selon la convention standard en cryptographie, PPT signifie temps polynomial probabiliste.}
\cmt{Le r\'educteur $\mathcal{B}$ est l'algorithme qui transforme une attaque en solution d'un probl\`eme difficile.}

\begin{table}[htbp]
\centering
\caption{Structure de l'Annexe \ref{chap:annexe_gamebased}}
\label{tab:ann_g_structure}
\small
\begin{tabular}{@{}llll@{}}
\toprule
\textbf{Section} & \textbf{Contenu} & \textbf{Propri\'et\'e}
& \textbf{Hypoth\`ese} \\
\midrule
\ref{sec:gamebased_primitives} & Primitives et d\'efinitions pr\'ealables & -- & -- \\
\ref{sec:gamebased_confidentialite} & Jeu IND-CONF + lien $\Conf$ & $\Conf$ & PRG \\
\ref{sec:gamebased_fiabilite} & Jeu BIND-REL + lien $\Relver$ & $\Relver$ & DLOG \\
\ref{sec:gamebased_opposabilite} & Jeu EXT-OPP + lien $\Opp$ & $\Opp$ & Fano \\
\ref{sec:gamebased_SOTP-Low} & R\'eductions SOTP-Low & $\Conf_{\mathrm{calc}}$, $\Relver$ & PRG, DLOG \\
\ref{sec:gamebased_SOTP-base} & Dimensionnement SOTP-Base & -- & JL \\
\ref{sec:gamebased_SOTP-High} & R\'eductions SOTP-High & $\Conf$, $\Relver$ & Groth16 \\
\ref{sec:gamebased_produitCRO} & \'enonc\'es complets + produit CLO & $\Conf_{\mathrm{calc}}\cdot\Relver\cdot\Opp$ & Tous \\
\ref{sec:gamebased_limitations} & Limitations et port\'ee & -- & -- \\
\bottomrule
\end{tabular}
\end{table}
\end{sloppypar}

% ============================================================================
\section{PRIMITIVES CRYPTOGRAPHIQUES ET D\'EFINITIONS PR\'EALABLES}
\addcontentsline{toc}{section}{Primitives et d\'efinitions pr\'ealables}
\label{sec:gamebased_primitives}
% ============================================================================
\begin{sloppypar}

\begin{cmtEtape}{Objectif de cette section.}
Rappeler les quatre primitives cryptographiques utilis\'ees dans la
famille SOTP, \`a savoir le PRG (couche Gold), l'engagement de Pedersen (couche Iron),
le zk-SNARK Groth16 (SOTP-High) et la transformation Johnson-Lindenstrauss
(SOTP-Base), ainsi que deux outils th\'eoriques, l'in\'egalit\'e de Fano et le DPI.
\end{cmtEtape}

\subsection{G\'en\'erateur pseudo-al\'eatoire (PRG)}
\noindent\textbf{Intuition.} Un PRG prend une graine courte
$s \in \{0,1\}^\lambda$ et produit une longue s\'equence
$G(s) \in \{0,1\}^{L(\lambda)}$, $L > \lambda$, indiscernable
d'une s\'equence vraiment al\'eatoire. Il fonde la couche
Gold de SOTP-Low, dans laquelle le pad $G(s)$ masque le message de sorte que, sans $s$,
aucun adversaire PPT ne peut r\'ecup\'erer $m$ depuis $m \oplus G(s)$.
\cmt{Le PRG \'etire une petite graine en un long pad pseudo-al\'eatoire.}
\cmt{Sans la graine, le pad est indiscernable d'un al\'ea v\'eritable.}

\begin{definition}[Sécurité PRG~\cite{GoldreichKrawczyk99}]\label{def:ann_prg}
Une fonction calculable en temps polynomial
$G\colon \{0,1\}^\lambda \to \{0,1\}^{L(\lambda)}$ avec $L(\lambda) > \lambda$
est un PRG sûr si, pour tout adversaire PPT $\mathcal{A}$, 
\begin{equation}
\mathsf{Adv}^{\mathsf{PRG}}_{\mathcal{A}}(\lambda) =
\left|
\Pr_{s \leftarrow \{0,1\}^\lambda}\!\big[\mathcal{A}(G(s)) = 1\big]
- \Pr_{y \leftarrow \{0,1\}^{L(\lambda)}}\!\big[\mathcal{A}(y) = 1\big]
\right| \leq \negl(\lambda)
\end{equation}
\cmtFormule{L'avantage mesure la capacit\'e \`a distinguer $G(s)$ d'un al\'ea pur.}
\cmt{L'avantage doit \^etre n\'egligeable pour tout adversaire PPT.}
\end{definition}

\noindent\textbf{Impl\'ementations pratiques.} SOTP-Low utilise
ChaCha20/AES-CTR, deux PRG standardis\'es atteignant cette d\'efinition.
\cmt{ChaCha20 et AES-CTR sont des PRG concrets, utilis\'es en pratique.}

\subsection{Engagement de Pedersen}

\noindent\textbf{Intuition.} L'engagement de Pedersen permet \`a un
prouveur de \guillemotleft~s'engager~\guillemotright sur $m$ sans le r\'ev\'eler (\emph{perfect hiding}),
tout en \'etant incapable de changer d'avis (\emph{computational binding}).
La construction $\mathsf{Com}(m, r) = g^m h^r$ exploite l'alg\`ebre des
groupes cycliques, de sorte que toute collision d'engagements permettrait de calculer
$\log_g h$, r\'esolvant le logarithme discret.
\cmt{Le hiding parfait signifie que le message est cach\'e inconditionnellement.}
\cmt{Par le binding computationnel, ouvrir sur deux messages casserait DLOG.}

\begin{definition}[Engagement~de~Pedersen~\cite{Pedersen91}]\label{def:ann_pedersen}
Soit $\mathbb{G}$ un groupe cyclique d'ordre premier $q$,
$g, h \in \mathbb{G}$ deux générateurs tels que $\log_g h$ est inconnu.
\begin{itemize}
\item $\mathsf{Com}(m, r) = g^m h^r \in \mathbb{G}$, $\quad m,r \in \mathbb{Z}_q$ (op\'eration de groupe)
\item $\mathsf{Open}(\sigma_2, m, r)$ accepte si $\sigma_2 = g^m h^r$
\end{itemize}
\textbf{Perfect Hiding.} Pour tout $m_0, m_1$, les engagements
$\mathsf{Com}(m_0, r_0)$ et $\mathsf{Com}(m_1, r_1)$ sont identiquement
distribués pour $r_0, r_1 \xleftarrow{\$} \mathbb{Z}_q$.
\cmt{Le masque $h^r$ est uniforme, ind\'ependant du message.}
\textbf{Computational Binding.} Aucun adversaire PPT ne peut trouver
$(m_0, r_0, m_1, r_1)$ avec $m_0 \neq m_1$ tel que
$g^{m_0} h^{r_0} = g^{m_1} h^{r_1}$ dans $\mathbb{G}$.
\cmt{Deux ouvertures distinctes r\'ev\`elent $\log_g h$.}
\end{definition}

\noindent\textbf{R\^ole dans SOTP.} La couche Iron de toute variante SOTP
est un engagement de Pedersen (sur la graine $s \in \mathbb{Z}_q$ pour SOTP-Low et
SOTP-Base, sur $m$ pour SOTP-High). Il garantit la fiabilit\'e de v\'erification
$\Relver$, puisqu'une preuve ne peut pas \^etre \guillemotleft~ouverte~\guillemotright sur deux messages
distincts.
\cmt{C'est la source de la fiabilit\'e de v\'erification dans toutes les variantes SOTP.}

\subsection{zk-SNARK Groth16}

\noindent\textbf{Intuition.} Un zk-SNARK pour une relation $\mathcal{R}$
permet de prouver qu'on conna\^it un t\'emoin $w$ tel que $(x,w) \in
\mathcal{R}$, sans r\'ev\'eler $w$, en une preuve de taille constante.
Une preuve Groth16 compte trois \'el\'ements de groupe, soit 128 octets (points compress\'es) sur BN254
ou 192 octets sur BLS12-381~; elle se v\'erifie en $O(1)$ couplages.
\cmt{Preuve de taille constante quel que soit le circuit.}
\cmt{La v\'erification se fait en temps constant via des couplages.}

\begin{definition}[zk-SNARK~Groth16~\cite{Groth16}]\label{def:ann_groth16}
$(\mathsf{Setup}, \mathsf{Prove}, \mathsf{Verify})$ satisfait les propri\'et\'es suivantes.

\noindent\textbf{Complétude parfaite.}
$\Pr[\mathsf{Verify}(\sigma_v,x,\pi)=1 \;|\; (x,w)\in\mathcal{R}] = 1$
\cmt{Une preuve honn\^ete est toujours accept\'ee.}

\noindent\textbf{Solidité connaissance.} Pour tout PPT $\mathcal{A}$,
il existe un extracteur $\mathcal{E}$ tel que
\begin{equation}
\Pr\!\left[
\mathsf{Verify}(\sigma_v,x,\pi)=1 \land (x,w)\notin\mathcal{R} :
(x,\pi)\leftarrow\mathcal{A}(\sigma_p),\; w\leftarrow\mathcal{E}(\sigma_p)
\right] \leq \negl(\lambda)
\end{equation}
\cmtFormule{L'extracteur r\'ecup\`ere un t\'emoin valide si la preuve passe.}
\cmt{Cette propri\'et\'e est plus forte que la soundness standard puisqu'on peut extraire le t\'emoin.}

\noindent\textbf{Zéro-connaissance parfaite.}
$\exists$ simulateur $\mathcal{S}$ tel que pour tout $(x,w)\in\mathcal{R}$,
\begin{equation}
\{\pi : \pi \leftarrow \mathsf{Prove}(\sigma_p,x,w)\}
\equiv
\{\pi : \pi \leftarrow \mathcal{S}(\sigma_p,x)\}
\quad\text{(distributions identiques)}
\end{equation}
\cmtFormule{Les vraies preuves et les preuves simul\'ees sont statistiquement identiques.}
\cmt{Cette \'egalit\'e exacte explique $\Conf \geq 1 - \negl$.}
\end{definition}

\noindent\textbf{Lecture pour SOTP-High.} La z\'ero-connaissance
\emph{parfaite}, jointe au masquage parfait de l'engagement $c_m$, implique que
$(\pi, c_m)$ est ind\'ependant de $m$ sachant $c$, d'o\`u $\Conf = 1$ au sens statistique.
La solidit\'e connaissance implique la fiabilit\'e de v\'erification
$\Relver \geq 1 - \negl(\lambda)$~; elle ne dit rien de la fiabilit\'e de
reconstruction, qui vaut $2^{-\lambda}$ lorsque $\Conf = 1$
(Remarque~\ref{rem:rel_ver}).
\cmt{La ZK parfaite donne la confidentialit\'e~; la solidit\'e donne la fiabilit\'e de v\'erification.}

\subsection{Transformation Johnson-Lindenstrauss (JL)}

\noindent\textbf{Intuition.} La JL-transform projette $N$ points de
$\mathbb{R}^n$ dans $\mathbb{R}^d$ ($d \ll n$) en pr\'eservant
approximativement toutes les distances par paires. Dans SOTP-Base,
les contextes $c \in \mathbb{R}^n$ sont compress\'es en $\Phi c \in
\mathbb{R}^d$, conservant la structure permettant l'extraction
institutionnelle.
\cmt{On compresse un espace de grande dimension en pr\'eservant les distances.}
\cmt{Cela permet de r\'eduire la taille de la preuve sans perdre l'information utile.}

\begin{definition}[Transformation~JL~\cite{JohnsonLindenstrauss84,DasguptaGupta03}]\label{def:ann_jl}
Pour $0 < \varepsilon < 1$, $N > 0$, $S \subset \mathbb{R}^n$ avec
$|S| = N$, il existe $\Phi\colon \mathbb{R}^n \to \mathbb{R}^d$ avec
$d = O(\varepsilon^{-2} \log N)$ telle que, pour tout $x,y \in S$,
\begin{equation}
(1-\varepsilon)\|x-y\|^2 \leq \|\Phi x-\Phi y\|^2 \leq (1+\varepsilon)\|x-y\|^2
\end{equation}
\cmtFormule{Les distances sont pr\'eserv\'ees \`a $\varepsilon$ pr\`es.}
\cmt{$d$ d\'epend logarithmiquement de $N$, pas lin\'eairement.}
\end{definition}

\noindent\textbf{Lemme de dimension pour SOTP-Base.}
La dimension minimale $d$ pour pr\'eserver $\Delta$ bits d'extraction
avec tol\'erance $\varepsilon^2$ est donn\'ee par la condition
\begin{equation}
d \geq \frac{4\Delta H(\mathcal{C})}{\varepsilon^2}
\label{eq:ann_jl_dim}
\end{equation}
\cmtFormule{La dimension cro\^it lin\'eairement en $\Delta H(C)$.}
\cmt{C'est le co\^ut de la compression pr\'eservant l'extraction.}

\subsection{Outils th\'eoriques fondamentaux}

\begin{lemma}[Inégalité~de~Fano~\cite{CoverThomas06}]\label{lem:ann_g_fano}
Soit $X$ une variable à $k$ valeurs, $\hat{X}$ un estimateur,
$p_e = \Pr[X \neq \hat{X}]$. Alors
\begin{equation}
H(X|\hat{X}) \leq h(p_e) + p_e \log_2(k-1)
\end{equation}
\cmtFormule{L'incertitude r\'esiduelle est born\'ee par la probabilit\'e d'erreur.}
\end{lemma}

\noindent\textbf{Lecture.} Faible erreur d'extraction $p_e$
$\Rightarrow$ faible incertitude r\'esiduelle $H(X|\hat{X})$
$\Rightarrow$ grande information mutuelle extraite. Ce lemme est le
pont direct entre le jeu EXT-OPP et la m\'etrique $\Opp$.
\cmt{Plus l'interpr\`ete se trompe peu, plus il extrait d'information.}
\cmt{C'est le lien formel entre le jeu EXT-OPP et la m\'etrique $\Opp$.}

\begin{lemma}[Traitement~de~l'Information~(DPI)~\cite{CoverThomas06}]\label{lem:ann_dpi}
Si $X \to Y \to Z$ forme une chaîne de Markov, alors
$I(X;Z) \leq I(X;Y)$.
\cmtFormule{Traiter une donn\'ee ne peut qu'\'eliminer de l'information.}
\end{lemma}

\noindent\textbf{Lecture.} Toute transformation de $Y$ ne peut
qu'augmenter l'incertitude sur $X$. Un adversaire observant $\sigma$
ne peut extraire au plus que $I(M;\Sigma|C)$ bits sur $m$, fondation
du lien entre $\mathsf{Adv}^{\mathsf{IND\text{-}CONF}}$ et $I(M;\Sigma|C)$.
\cmt{Le DPI justifie que la meilleure attaque est born\'ee par $I(M;\Sigma|C)$.}
\end{sloppypar}

% ============================================================================
\section{JEU IND-CONF~: CONFIDENTIALIT\'E GAME-BASED}
\addcontentsline{toc}{section}{Jeu IND-CONF}
\label{sec:gamebased_confidentialite}
% ============================================================================
\begin{sloppypar}

\begin{cmtEtape}{Objectif de cette section.}
D\'efinir le jeu IND-CONF qui capture la confidentialit\'e, puis
\'etablir son lien avec la m\'etrique $\Conf$ du cadre CLO.
\end{cmtEtape}

\subsection{D\'efinition du jeu}

\noindent\textbf{Intuition.} L'adversaire choisit deux messages
$m_0, m_1$, re\c{c}oit la preuve de l'un des deux (au choix secret du
challenger) puis doit deviner lequel. Si le protocole est confidentiel,
l'adversaire ne fait pas mieux qu'une devinette al\'eatoire
(avantage $\approx 0$).
\cmt{L'adversaire ne doit pas pouvoir distinguer les preuves des deux messages.}

\begin{definition}[Jeu~IND-CONF (section~\ref{sec:sotp_construction})]\label{def:ann_ind_conf}
Pour un protocole $\mathcal{P} = (\mathsf{Setup},\mathsf{Prove},
\mathsf{Verify})$ et un adversaire $\mathcal{A}$, le jeu se d\'eroule comme suit.
\begin{enumerate}
\item $\mathcal{A}$ reçoit $pp \leftarrow \mathsf{Setup}(1^\lambda)$
\item $\mathcal{A}$ choisit $m_0,m_1 \in \mathcal{M}_\lambda$
      ($|m_0|=|m_1|$) et $c \in \mathcal{C}_\lambda$
      \hfill\textit{(phase de défi)}
      \cmt{L'adversaire choisit les deux messages \`a distinguer.}
\item Le challenger tire $b \xleftarrow{\$} \{0,1\}$ et calcule
      $\sigma \leftarrow \mathsf{Prove}(pp, m_b, c)$
      \cmt{Le challenger chiffre l'un des deux messages au hasard.}
\item $\mathcal{A}$ reçoit $\sigma$ et émet $b' \in \{0,1\}$
      \hfill\textit{(phase de réponse)}
      \cmt{L'adversaire devine lequel des deux messages a \'et\'e chiffr\'e.}
\item L'expérience produit 1 si $b' = b$, sinon 0
\end{enumerate}
\begin{equation}
\mathsf{Adv}^{\mathsf{IND\text{-}CONF}}_{\mathcal{A},\mathcal{P}}(\lambda)
= \left|\Pr\!\left[
\mathsf{Exp}^{\mathsf{IND\text{-}CONF}}_{\mathcal{A},\mathcal{P}}(\lambda)
= 1\right] - \frac{1}{2}\right|
\end{equation}
\cmtFormule{L'avantage est l'\'ecart \`a la devinette al\'eatoire (probabilit\'e 1/2).}
\cmt{Un avantage n\'egligeable signifie confidentialit\'e.}
\end{definition}

\subsection{Lien avec $\Conf$ au sens de la th\'eorie de l'information}

\begin{remark}[Pont IND-CONF $\leftrightarrow$ $\Conf$]
\label{rem:ann_ind_conf_link}
Pour tout distingueur $\mathcal{A}$, m\^eme non born\'e en calcul, on a
\begin{equation}
\mathsf{Adv}^{\mathsf{IND\text{-}CONF}}_{\mathcal{A},\mathcal{P}}(\lambda)
\;\leq\; \frac{1 - \Conf(\mathcal{P})}{2},
\end{equation}
\cmtFormule{L'avantage IND-CONF est born\'e par la distinguabilit\'e statistique.}
avec \'egalit\'e pour le test de Neyman--Pearson sur la paire la plus
distinguable (Remarque~\ref{rem:jeu_ind_conf_link}). Un protocole avec
$\Conf \geq 1 - \negl(\lambda)$ r\'esiste donc au jeu IND-CONF. La r\'eciproque
est fausse, car la r\'esistance \`a IND-CONF contre les adversaires PPT est une
propri\'et\'e calculatoire, compatible avec une confidentialit\'e statistique faible.
Les r\'eductions de cette annexe \'etablissent des garanties calculatoires, que
nous notons $\Conf_{\mathrm{calc}}$ lorsqu'il faut les distinguer de $\Conf$.
\cmt{La confidentialit\'e statistique implique la confidentialit\'e calculatoire, pas l'inverse.}
\end{remark}

\end{sloppypar}

% ============================================================================
\section{JEU BIND-REL~: FIABILIT\'E GAME-BASED}
\addcontentsline{toc}{section}{Jeu BIND-REL}
\label{sec:gamebased_fiabilite}
% ============================================================================
\begin{sloppypar}

\begin{cmtEtape}{Objectif de cette section.}
D\'efinir le jeu BIND-REL qui capture la fiabilit\'e (non-\'equivocit\'e),
puis \'etablir son lien avec la m\'etrique $\Relver$.
\end{cmtEtape}

\subsection{D\'efinition du jeu}

\noindent\textbf{Intuition.} Le jeu capture la \emph{non-\'equivocit\'e},
selon laquelle un adversaire ne doit pas pouvoir produire une preuve $\sigma$ ouvrable
sur deux messages distincts $m_0 \neq m_1$. Contrairement \`a IND-CONF,
l'adversaire choisit librement $(c, \sigma, s_0, s_1)$, mod\'elisant
un agent qui forge d\'elib\'er\'ement une preuve ambiguë.
\cmt{L'adversaire essaie d'ouvrir une m\^eme preuve sur deux messages diff\'erents.}

\begin{definition}[Jeu~BIND-REL (section~\ref{sec:sotp_construction})]\label{def:ann_bind_rel}
Pour un protocole $\mathcal{P}$ avec $\mathsf{Open}$ et un adversaire
$\mathcal{A}$, le jeu se d\'eroule comme suit.
\begin{enumerate}
\item $\mathcal{A}$ reçoit $pp \leftarrow \mathsf{Setup}(1^\lambda)$
\item $\mathcal{A}$ produit $(c, \sigma, s_0, s_1)$
      \cmt{L'adversaire forge une preuve et deux t\'emoins.}
\item Calculer $m_i = \mathsf{Open}(pp, \sigma, s_i)$ pour $i \in \{0,1\}$
\item $\mathcal{A}$ \emph{gagne} si et seulement si
  $\mathsf{Verify}(pp,\sigma,c)=1$,
  $m_0 \neq m_1$,
  $m_0 \neq \bot$ et $m_1 \neq \bot$
  \cmt{Il gagne si la preuve passe et les deux messages sont distincts et valides.}
\end{enumerate}
\begin{equation}
\mathsf{Win}^{\mathsf{BIND\text{-}REL}}_{\mathcal{A},\mathcal{P}}(\lambda)
= \Pr\!\left[\mathcal{A} \text{ gagne dans }
\mathsf{Exp}^{\mathsf{BIND\text{-}REL}}_{\mathcal{A},\mathcal{P}}(\lambda)
\right]
\end{equation}
\cmtFormule{Probabilit\'e que l'adversaire produise une collision d'engagement.}
\end{definition}

\subsection{Lien avec $\Relver$ (fiabilit\'e de v\'erification)}

\begin{remark}[Pont BIND-REL $\leftrightarrow$ $\Relver$]
\label{rem:ann_bind_rel_link}
Pour un protocole complet dont la probabilit\'e de gain BIND-REL est au plus
$\varepsilon_{\mathsf{bind}}$, on a
\begin{equation}
\Relver \geq 1 - \varepsilon_{\mathsf{bind}} - \negl(\lambda)
\end{equation}
\cmtFormule{Fiabilit\'e de v\'erification born\'ee inf\'erieurement par $1 - \varepsilon_{\mathsf{bind}}$.}
(Remarque~\ref{rem:jeu_bind_rel_link}). BIND-REL ne dit rien de la fiabilit\'e de
reconstruction $\Relrec$, puisqu'un engagement parfaitement masquant et liant a
$\Relrec = 2^{-\lambda}$.
\cmt{Un engagement liant ne r\'ev\`ele pas pour autant le message.}
\end{remark}
\end{sloppypar}

% ============================================================================
\section{JEU EXT-OPP~: OPPOSABILIT\'E GAME-BASED}
\addcontentsline{toc}{section}{Jeu EXT-OPP}
\label{sec:gamebased_opposabilite}
% ============================================================================
\begin{sloppypar}

\begin{cmtEtape}{Objectif de cette section.}
D\'efinir le jeu EXT-OPP qui capture l'opposabilit\'e institutionnelle,
puis \'etablir son lien avec la m\'etrique $\Opp$ via l'in\'egalit\'e de Fano.
\end{cmtEtape}

\subsection{D\'efinition du jeu}

\noindent\textbf{Intuition.} Le jeu mod\'elise la capacit\'e d'un interpr\`ete
institutionnel \`a extraire le bon verdict depuis une preuve. L'\guillemotleft~adversaire~\guillemotright
ici est l'interpr\`ete lui-m\^eme (juge, r\'egulateur), dont on mesure la probabilit\'e
d'\emph{erreur}. Le verdict de r\'ef\'erence $v^*$ est celui que l'institution rendrait en
pleine connaissance de cause (D\'efinition~\ref{def:interp_admissible})~;
l'institution doit rendre une d\'ecision align\'ee sur le contexte l\'egal, au-del\`a
de la seule v\'erification de la preuve.
\cmt{Contrairement aux autres jeux, l'adversaire est ici l'interpr\`ete, qui n'est pas malveillant.}
\cmt{On mesure sa probabilit\'e de se tromper sur le verdict.}

\begin{definition}[Jeu~EXT-OPP (section~\ref{sec:sotp_construction})]\label{def:ann_ext_opp}
Pour un protocole $\mathcal{P}$, un interprète $\mathsf{Interpret}$
de capacité $\Delta$ et une distribution $\mathcal{D}_C$ fixée à l'avance, le jeu se d\'eroule comme suit.
\begin{enumerate}
\item Le challenger tire $c \leftarrow \mathcal{D}_C$,
      $m \leftarrow \mathcal{M}_\lambda$
      \cmt{Le contexte est tir\'e selon la distribution du domaine.}
\item Le challenger calcule $\sigma \leftarrow \mathsf{Prove}(pp,m,c)$
\item L'interprète reçoit $(\sigma,c)$ et produit
      $v \leftarrow \mathsf{Interpret}(pp,\sigma,c,\Delta)$
      \cmt{L'interpr\`ete rend un verdict \`a partir de la preuve et du contexte.}
\item Soit $v^*$ le verdict de r\'ef\'erence, r\'ealisation de la variable $V$
      (qui peut d\'ependre du message, du contexte et des ressources
      institutionnelles, D\'efinition~\ref{def:interp_admissible})
      \cmt{Le verdict de r\'ef\'erence est celui que l'institution rendrait en pleine connaissance de cause.}
\item L'interprète \emph{gagne} si $v = v^*$
\end{enumerate}
\begin{equation}
\mathsf{Err}^{\mathsf{EXT\text{-}OPP}}_{\mathsf{Interpret},\mathcal{P},\mathcal{D}_C}(\lambda)
= \Pr\!\left[
\mathsf{Interpret} \text{ perd dans }
\mathsf{Exp}^{\mathsf{EXT\text{-}OPP}}\right]
\end{equation}
\cmtFormule{Probabilit\'e que l'interpr\`ete rende un verdict incorrect.}
\end{definition}

\noindent\textbf{Pourquoi $\mathcal{D}_C$ est fix\'ee \`a l'avance ?}
L'opposabilit\'e d\'epend du domaine institutionnel. Un protocole tr\`es
opposable dans un contexte contractuel simple peut l'\^etre peu dans un
contexte r\'eglementaire complexe. Fixer $\mathcal{D}_C$ param\'etrise le
jeu par le domaine.
\cmt{Le domaine institutionnel conditionne la difficult\'e de l'extraction.}

\subsection{Lien avec $\Opp$ via Fano}

\begin{remark}[Pont EXT-OPP $\leftrightarrow$ $\Opp$]
\label{rem:ann_ext_opp_link}
Par Fano (Lemme~\ref{lem:ann_g_fano}) appliqu\'e \`a l'estimation de $V$ \`a
partir de $(\Sigma,C)$, on obtient
\begin{equation}
\Opp = 1 - \frac{H(V|\Sigma,C)}{H(V|C)}
\geq 1 - \frac{h(\varepsilon_{\mathsf{ext}})
  + \varepsilon_{\mathsf{ext}}\log_2(|\mathcal{V}_\lambda|-1)}
  {H(V|C)}
\label{eq:ann_fano_opp}
\end{equation}
\cmtFormule{L'opposabilit\'e est born\'ee inf\'erieurement par $1 - \varepsilon_{\mathsf{ext}}$.}
où $\varepsilon_{\mathsf{ext}} = \mathsf{Err}^{\mathsf{EXT\text{-}OPP}}$
et $h(\cdot)$ est l'entropie binaire.
Un interprète avec faible erreur d'extraction implique une forte
opposabilité.
\cmt{Moins l'interpr\`ete se trompe, plus l'opposabilit\'e est grande.}
\end{remark}
\end{sloppypar}

% ============================================================================
\section{R\'EDUCTIONS FORMELLES POUR SOTP-LOW}
\addcontentsline{toc}{section}{R\'eductions SOTP-Low}
\label{sec:gamebased_SOTP-Low}
% ============================================================================
\begin{sloppypar}
SOTP-Low op\`ere dans le r\'egime \`a contexte rare ($H(\mathcal{C}) < H^*(\lambda)$).
Ses deux r\'eductions principales ciblent PRG (confidentialit\'e) et
logarithme discret (fiabilit\'e).

\begin{cmtEtape}{Objectif de cette section.}
\'Etablir les deux r\'eductions formelles qui relient SOTP-Low aux
hypoth\`eses standard PRG (couche Gold) et DLOG (couche Iron).
\end{cmtEtape}

\subsection{R\'eduction PRG pour la confidentialit\'e}

\noindent\textbf{Intuition.} La couche Gold chiffre $m$ par le pad $G(s)$.
Si $\mathcal{A}$ distingue les preuves de $m_0$ vs $m_1$, il distingue
n\'ecessairement $G(s)$ d'une cha\^ine al\'eatoire, contredisant la s\'ecurit\'e PRG.
\cmt{Toute attaque contre la confidentialit\'e devient une attaque contre le PRG.}

\begin{theorem*}[SOTP-Low~:~Réduction~PRG (section~\ref{sec:sotp_construction})]\label{thm:ann_sotp_low_conf}
Soit $\mathcal{A}$ adversaire PPT contre IND-CONF de SOTP-Low avec
avantage $\varepsilon$. Alors $\exists$ adversaire PPT $\mathcal{B}$
contre le PRG $G$ avec avantage $\geq \varepsilon/2 - \negl(\lambda)$.
\cmt{L'avantage contre le PRG est au moins la moiti\'e de l'avantage contre IND-CONF.}
\end{theorem*}

\begin{Preuve}

\begin{cmtEtape}{Construction du r\'educteur $\mathcal{B}$.}
$\mathcal{B}$ re\c{c}oit un d\'efi PRG (soit $G(s)$, soit uniforme) et
simule l'environnement IND-CONF pour $\mathcal{A}$. Il utilise son
d\'efi comme pad masquant, puis observe si $\mathcal{A}$ devine
correctement $b$.
\end{cmtEtape}

\textbf{Construction de $\mathcal{B}$.} $\mathcal{B}$ re\c{c}oit un d\'efi
PRG $y \in \{0,1\}^L$ (soit $G(s)$ soit uniforme) et simule le jeu
IND-CONF pour $\mathcal{A}$ de la mani\`ere suivante.
\begin{enumerate}
\item G\'en\'erer $(\mathbb{G},q,g,h)$ comme dans le Setup SOTP-Low
      \cmt{$\mathcal{B}$ peut g\'en\'erer les param\`etres publics car ils ne d\'ependent pas du PRG.}
\item Quand $\mathcal{A}$ soumet $(m_0,m_1,c)$, choisir
      $b \xleftarrow{\$} \{0,1\}$
      \cmt{$\mathcal{B}$ choisit un bit secret comme dans le vrai jeu.}
\item Poser $\sigma_1 = m_b \oplus y$ \hfill (d\'efi PRG comme pad)
      \cmt{Le d\'efi PRG joue le r\^ole du pad.}
\item Poser $\sigma_2 = g^{s'}h^\rho$ pour
      $s',\rho \xleftarrow{\$} \mathbb{Z}_q$ \hfill (Iron simul\'ee)
      \cmt{$\mathcal{B}$ simule la couche Iron sans conna\^itre $s$.}
\item Poser $\sigma_3 = (\phi(c), \mathsf{Hash}(c), V(c))$
      \hfill (Clay d\'eterministe)
      \cmt{La couche Clay est simul\'ee exactement comme dans SOTP-Low.}
\item Retourner $\sigma = (\sigma_1,\sigma_2,\sigma_3)$ \`a $\mathcal{A}$
\end{enumerate}
Si $\mathcal{A}$ \'emet $b'=b$, $\mathcal{B}$ retourne 1 ; sinon 0.
\cmt{$\mathcal{B}$ retourne 1 si $\mathcal{A}$ a devin\'e correctement.}

\textbf{Analyse.}
\begin{itemize}
\item \emph{Cas $y = G(s)$.} $\mathcal{A}$ voit une vraie preuve SOTP-Low,
      gagne avec probabilit\'e $1/2 + \varepsilon$.
      $\mathcal{B}$ retourne 1 avec probabilit\'e $1/2+\varepsilon$.
      \cmt{Dans ce cas, $\mathcal{B}$ a bien simul\'e le vrai jeu.}
\item \emph{Cas $y$ uniforme.} $\sigma_1 = m_b \oplus U_{|m|}$ est un
      one-time pad parfait, $\mathcal{A}$ n'a aucune information sur $b$.
      $\mathcal{B}$ retourne 1 avec probabilit\'e $1/2$.
      \cmt{Dans ce cas, $\mathcal{A}$ ne peut pas mieux faire que 1/2.}
\end{itemize}
Donc $\mathsf{Adv}^{\mathsf{PRG}}_{\mathcal{B}}(\lambda) \geq
\varepsilon/2 - \negl(\lambda)$.
\cmtFormule{Conclusion de la r\'eduction.}
\end{Preuve}

\subsection{R\'eduction logarithme discret pour la fiabilit\'e}

\noindent\textbf{Intuition.} La couche Iron est
$\sigma_2 = g^s h^\rho$. Pour gagner BIND-REL, l'adversaire doit trouver
deux ouvertures $s_0 \neq s_1$ pour le m\^eme $\sigma_2$, ce qui implique
$g^{s_0}h^{\rho_0} = g^{s_1}h^{\rho_1}$, permettant de calculer
$\log_g h$, r\'esolvant DLOG.
\cmt{Deux ouvertures distinctes donnent directement le logarithme discret.}

\begin{theorem}[SOTP-Low~:~Réduction~DLOG (section~\ref{sec:sotp_construction})]\label{thm:ann_sotp_low_rel}
On suppose la graine tir\'ee dans $\mathbb{Z}_q$ et l'engagement de la forme
$\sigma_2 = g^{s}h^{\rho}$ (Algorithme~\ref{alg:sotp_low}). Soit $\mathcal{A}$ PPT
gagnant BIND-REL pour SOTP-Low (proc\'edure d'ouverture $\mathsf{Open}$) avec
probabilit\'e $\varepsilon$. Alors $\exists$ adversaire PPT $\mathcal{B}$ r\'esolvant
DLOG dans $\mathbb{G}$ avec probabilit\'e $\geq \varepsilon$.
\cmt{L'avantage contre DLOG est le m\^eme que contre BIND-REL.}
\end{theorem}

\begin{Preuve}

\begin{cmtEtape}{Construction du r\'educteur DLOG.}
$\mathcal{B}$ re\c{c}oit une instance DLOG $(g, h)$ et utilise l'attaque
de $\mathcal{A}$ contre BIND-REL pour extraire $\log_g h$. La cl\'e est
que deux ouvertures valides sur la m\^eme preuve donnent une \'equation
alg\'ebrique r\'esoluble.
\end{cmtEtape}

Soit une instance DLOG $(g,h)$ avec $x = \log_g h$ inconnu.
\begin{enumerate}
\item $\mathcal{B}$ configure SOTP-Low avec $(g,h)$ donn\'es
      \cmt{$\mathcal{B}$ utilise l'instance DLOG comme param\`etres publics~; leur loi est celle du Setup.}
\item $\mathcal{A}$ produit $(c,\,\sigma{=}(\sigma_1,\sigma_2,\sigma_3))$ et deux
      ouvertures accept\'ees $(m_0, s_0, \rho_0)$, $(m_1, s_1, \rho_1)$ avec
      $m_0 \neq m_1$
      \cmt{$\mathcal{A}$ fournit sa collision d'ouverture.}
\item Comme $\sigma_1 = m_0 \oplus \mathsf{PRG}(s_0) = m_1 \oplus \mathsf{PRG}(s_1)$ et
      $m_0 \neq m_1$, on a $\mathsf{PRG}(s_0) \neq \mathsf{PRG}(s_1)$, donc
      $s_0 \neq s_1$ dans $\mathbb{Z}_q$
      \cmt{C'est ici que le tirage de la graine dans $\mathbb{Z}_q$ est indispensable, car avec une graine de $\{0,1\}^\lambda$ engag\'ee par $\mathsf{int}(s) \bmod q$, $s$ et $s+q$ donneraient le m\^eme engagement.}
\item Les deux ouvertures donnent $g^{s_0}h^{\rho_0} = g^{s_1}h^{\rho_1}$, donc
      $g^{s_0-s_1} = h^{\rho_1-\rho_0}$~; comme $s_0 \neq s_1$ et $g$ engendre
      $\mathbb{G}$, on a $\rho_1 \neq \rho_0$
      \cmtFormule{On isole $\log_g h = x$ dans l'\'equation.}
\item Calculer
\begin{equation}
x = \bigl(s_0-s_1\bigr)\cdot\bigl(\rho_1-\rho_0\bigr)^{-1} \bmod q
\end{equation}
\cmtFormule{On r\'esout l'\'equation lin\'eaire pour obtenir $x = \log_g h$.}
\end{enumerate}
$\mathcal{B}$ r\'eussit chaque fois que $\mathcal{A}$ produit une
attaque valide, avec probabilit\'e $\varepsilon$.
\cmt{La r\'eduction est exacte, sans perte d'avantage.}
\end{Preuve}

\noindent\textbf{Signification pratique.} Avec Curve25519 ($q \approx 2^{252}$,
s\'ecurit\'e 128 bits), casser la fiabilit\'e de SOTP-Low est aussi difficile
que r\'esoudre le DLP sur les courbes elliptiques.
\cmt{Le niveau de s\'ecurit\'e de SOTP-Low h\'erite de celui de DLOG sur la courbe.}
\end{sloppypar}

% ============================================================================
\section{R\'EDUCTION JL POUR SOTP-BASE}
\addcontentsline{toc}{section}{R\'eduction SOTP-Base}
\label{sec:gamebased_SOTP-base}
% ============================================================================
\begin{sloppypar}
SOTP-Base op\`ere dans le r\'egime de transition
($H^*(\lambda) \leq H(\mathcal{C}) \leq \lambda/\log_2\lambda$).
Sa couche Clay applique une projection JL $\Phi\colon \mathbb{R}^n \to
\mathbb{R}^d$ avec $d = \lceil 4\Delta H(\mathcal{C})/\varepsilon^2
\rceil$.

\begin{cmtEtape}{Objectif de cette section.}
Montrer que l'erreur d'extraction de l'interpr\`ete est directement
reli\'ee \`a la qualit\'e de la pr\'eservation JL. Si la projection
d\'eforme trop les distances, l'extraction devient erron\'ee.
\end{cmtEtape}

\noindent\textbf{Intuition.} Si $\Phi$ d\'eforme trop certaines distances,
les contextes voisins deviennent indiscernables dans l'espace compress\'e,
causant des erreurs d'extraction. La r\'eduction ci-dessous quantifie
pr\'ecis\'ement cette relation.
\cmt{La compression JL est le seul point de perte en SOTP-Base.}

\begin{theorem}[SOTP-Base~:~Dimensionnement~JL (section~\ref{sec:sotp_construction})]\label{thm:ann_sotp_base_opp}
Soit $\mathcal{C} = \{c_1,\ldots,c_N\}$ l'ensemble des contextes plong\'es, $N =
2^{H(\mathcal{C})}$ et $\delta, \varepsilon' \in (0,1)$. Si la dimension de
projection v\'erifie
\begin{equation}
d \;\geq\; \frac{8}{\delta^2}\,\ln\!\left(\frac{N^2}{\varepsilon'}\right)
= O\!\left(\delta^{-2}\bigl(H(\mathcal{C}) + \log(1/\varepsilon')\bigr)\right),
\end{equation}
alors, avec probabilit\'e au moins $1 - \varepsilon'$ sur le tirage de $\Phi$,
$\Phi$ pr\'eserve toutes les distances par paires de $\mathcal{C}$ \`a un facteur
$(1 \pm \delta)$ pr\`es.
\cmt{La pr\'eservation JL conditionne l'extraction~; la r\'eciproque est fausse.}
\end{theorem}

\begin{Preuve}
Par le lemme JL, pour chaque paire $(i,j)$,
$\Pr\bigl[\,\bigl|\|\Phi c_i-\Phi c_j\|^2/\|c_i-c_j\|^2 - 1\bigr| > \delta\,\bigr]
\leq 2e^{-d\delta^2/8}$~; une borne de l'union sur les $N^2/2$ paires donne une
probabilit\'e d'\'echec au plus $N^2 e^{-d\delta^2/8}$, inf\'erieure \`a $\varepsilon'$
d\`es que $d \geq (8/\delta^2)\ln(N^2/\varepsilon')$.
\cmtFormule{Borne de l'union sur les paires de contextes.}
La pr\'eservation des distances est une condition \emph{n\'ecessaire} pour qu'une
fonction d'extraction fond\'ee sur les distances distingue les contextes~; elle
n'est pas suffisante pour garantir l'information extractible sur le verdict.
Elle ne fournit donc aucune minoration de $\Opp$. Cette dimension, en
$O(H(\mathcal{C}))$, est distincte de la r\`egle de dimensionnement
$\lceil 4\Delta H(\mathcal{C})/\varepsilon^2 \rceil$ du texte principal, qui est
un choix de conception.
\end{Preuve}

\noindent\textbf{Interpr\'etation.} SOTP-Base choisit la dimension
$d$ garantissant la pr\'eservation des distances avec la probabilit\'e voulue~;
le lien avec l'extraction du verdict reste une hypoth\`ese de conception.
\cmt{Le choix de $d$ est un compromis entre taille de preuve et qualit\'e d'extraction.}
\end{sloppypar}

% ============================================================================
\section{R\'EDUCTIONS POUR SOTP-HIGH}
\addcontentsline{toc}{section}{R\'eductions SOTP-High}
\label{sec:gamebased_SOTP-High}
% ============================================================================
\begin{sloppypar}
SOTP-High op\`ere dans le r\'egime \`a contexte riche
($H(\mathcal{C}) > \lambda/\log_2\lambda$) et s'appuie sur Groth16
pour sa couche Gold. Sa couche Clay encode $h_c = \mathsf{Poseidon}(c)$.

\begin{cmtEtape}{Objectif de cette section.}
Montrer que $\Conf$ et $\Relver$ de SOTP-High h\'eritent des propri\'et\'es de
Groth16 (z\'ero-connaissance parfaite, solidit\'e connaissance)~; l'opposabilit\'e
n'est pas couverte par une r\'eduction.
\end{cmtEtape}

\begin{theorem}[SOTP-High~:~Sécurité~héritée~de~Groth16 (section~\ref{sec:sotp_construction})]\label{thm:ann_sotp_high}
Pour la relation $\mathcal{R} = \{((h_c, c_m), (m, c, \rho)) : h_c = \mathsf{Poseidon}(c)
\wedge c_m = g^m h^\rho \wedge \mathsf{Conforme}(m,c) = 1\}$, les deux propri\'et\'es suivantes sont v\'erifi\'ees.
\begin{enumerate}
\item \textbf{Confidentialité.} La z\'ero-connaissance parfaite
      de Groth16 et le masquage parfait de $c_m$ donnent $\Conf = 1$.
      \cmt{La preuve et l'engagement sont ind\'ependants de $m$ sachant $c$.}
\item \textbf{Fiabilité de v\'erification.} La solidité connaissance de
      Groth16 donne $\Relver \geq 1 - \negl(\lambda)$.
      \cmt{Un prouveur qui fait accepter conna\^it un t\'emoin valide.}
\end{enumerate}
\end{theorem}

\begin{Preuve}

\textbf{Confidentialit\'e.} Par la ZK parfaite, les preuves honn\^etes et les
preuves simul\'ees $\mathcal{S}(\sigma_p, (h_c, c_m))$ ont la m\^eme loi. Le
simulateur n'utilise que les entr\'ees publiques~; $h_c$ ne d\'epend que de $c$ tandis que
$c_m$ est uniforme dans $\mathbb{G}$ quel que soit $m$. La loi de
$\Sigma = (h_c, \pi, c_m, \mathsf{metadata})$ sachant $(M = m, C = c)$ ne d\'epend donc
pas de $m$, d'o\`u $\Conf = 1$.
\cmtFormule{La ZK parfaite implique l'indistinguabilit\'e inconditionnelle.}

\textbf{Fiabilit\'e de v\'erification.} Pour tout $\mathcal{A}$ PPT produisant une
preuve acceptante $\pi$, l'extracteur $\mathcal{E}$ de Groth16 (qui dispose du
code de $\mathcal{A}$ et exploite l'hypoth\`ese q-PKE~; ce n'est pas le
v\'erificateur) produit un t\'emoin $(m,c,\rho)$ avec $((h_c,c_m),(m,c,\rho)) \in
\mathcal{R}$, sauf avec probabilit\'e n\'egligeable. Ainsi
\begin{equation}
\Pr\!\left[\mathsf{Verify}=1 \land (m,c,\rho)\notin\mathcal{R}\right]
\leq \mathsf{Adv}^{\mathsf{sound}}(\lambda) \leq \negl(\lambda)
\;\Rightarrow\; \Relver \geq 1-\negl(\lambda).
\end{equation}
\cmtFormule{La solidit\'e connaissance implique l'absence de forge.}

\textbf{Opposabilit\'e.} Aucune borne n'est d\'eriv\'ee ici, car $h_c$ et les
m\'etadonn\'ees sont des fonctions du contexte public, donc $I(V;h_c|C) = 0$~; ce que
la preuve apporte au verdict tient \`a l'assertion certifi\'ee
$\mathsf{Conforme}(m,c) = 1$ (Remarque~\ref{rem:heuristique_sotp_high}).
\end{Preuve}

\begin{remark}[La ZK parfaite] La ZK \emph{parfaite}
de Groth16 signifie que m\^eme un adversaire \`a calcul infini ne peut
distinguer vraies preuves de simul\'ees. C'est une garantie inconditionnelle
une fois le CRS fix\'e, propri\'et\'e forte, qui fait de SOTP-High
la variante la plus protectrice pour les contextes \`a haute entropie.
\cmt{La ZK parfaite est inconditionnelle.}
\end{remark}
\end{sloppypar}

% ============================================================================
\section{\'ENONC\'ES COMPLETS ET PRODUIT CRO GAME-BASED}
\addcontentsline{toc}{section}{\'Enonc\'es complets et produit CLO}
\label{sec:gamebased_produitCRO}
% ============================================================================
\begin{sloppypar}

\begin{cmtEtape}{Objectif de cette section.}
Rassembler les bornes game-based des trois jeux et en d\'eduire une
minoration du produit CLO \`a l'\'echelle op\'erationnelle. Elle ne se compare pas
directement au Trilemme CRO (Remarque~\ref{rem:ann_pas_de_contradiction}).
\end{cmtEtape}

\subsection{S\'ecurit\'e compl\`ete de SOTP-Low}

\begin{theorem}[SOTP-Low~:~Bornes~game-based~compl\`etes (section~\ref{sec:sotp_construction})]\label{thm:ann_sotp_low_games}
Sous PRG et DLOG dans $\mathbb{G}$ d'ordre $q$, on a
\begin{align}
\mathsf{Adv}^{\mathsf{IND\text{-}CONF}}_{\mathcal{A},\mathsf{SOTP\text{-}Low}}(\lambda)
&\leq 2\cdot\mathsf{Adv}^{\mathsf{PRG}}_{\mathcal{B}}(\lambda)
    + \frac{1}{q}
\label{eq:ann_sotp_low_conf} \\[4pt]
\mathsf{Win}^{\mathsf{BIND\text{-}REL}}_{\mathcal{A},\mathsf{SOTP\text{-}Low}}(\lambda)
&\leq \mathsf{Adv}^{\mathsf{DLOG}}_{\mathcal{B}'}(\lambda)
    + \frac{1}{q}
\label{eq:ann_sotp_low_rel}
\end{align}
\cmtFormule{Confidentialit\'e calculatoire born\'ee par l'avantage PRG.}
\cmtFormule{Fiabilit\'e de v\'erification born\'ee par l'avantage DLOG.}
pour tout PPT $\mathcal{A}$.
\end{theorem}

\begin{remark}[Opposabilit\'e de SOTP-Low~: pas de borne game-based]
\label{rem:ann_sotp_low_opp_statut}
La formulation initiale de ce th\'eor\`eme contenait une troisi\`eme in\'egalit\'e,
$\mathsf{Err}^{\mathsf{EXT\text{-}OPP}} \leq 1 - \Delta/(\alpha H(\mathcal{C}))$,
tir\'ee du fait que la couche Clay $\sigma_3 = (\phi(c), \mathsf{Hash}(c), V(c))$ d\'etermine
le verdict. La pr\'esente formulation ne la reprend pas, car $\sigma_3$ est une fonction du
contexte public $C$, donc $I(V;\sigma_3|C) = 0$~; elle ne contribue ainsi en rien \`a
l'opposabilit\'e au sens de la D\'efinition~\ref{def:interp_admissible}. Ce que
SOTP-Low apporte au verdict passe par l'ouverture $(m,s,\rho)$ fournie \`a
l'institution, ce qui rel\`eve de l'heuristique de conception
(Remarque~\ref{rem:facteur_limitant_low}), sans r\'eduction associ\'ee.
\end{remark}

\subsection{Tableau r\'ecapitulatif des r\'eductions}

\begin{table}[htbp]
\centering
\caption{Synth\`ese des r\'eductions game-based pour la famille SOTP}
\label{tab:ann_g_reductions}
\begin{tabular}{@{}lllll@{}}
\toprule
\textbf{Variante} & \textbf{Propri\'et\'e} & \textbf{Jeu} &
\textbf{Hypoth\`ese} & \textbf{Borne} \\
\midrule
SOTP-Low & $\Conf_{\mathrm{calc}}$ & IND-CONF & PRG
  & $\mathsf{Adv} \leq 2\varepsilon_{\mathsf{PRG}} + 1/q$ \\
SOTP-Low & $\Relver$ & BIND-REL & DLOG
  & $\mathsf{Win} \leq \varepsilon_{\mathsf{DLOG}} + 1/q$ \\
SOTP-Low & $\Opp$ & -- & --
  & aucune (Rem.~\ref{rem:ann_sotp_low_opp_statut}) \\
\midrule
SOTP-Base & $\Conf_{\mathrm{calc}}$ & IND-CONF & PRG
  & $\mathsf{Adv} \leq 2\varepsilon_{\mathsf{PRG}} + \negl$ \\
SOTP-Base & $\Relver$ & BIND-REL & DLOG
  & $\mathsf{Win} \leq \varepsilon_{\mathsf{DLOG}} + \negl$ \\
SOTP-Base & -- & -- & Lemme JL
  & dimensionnement de $d$ (Th.~\ref{thm:ann_sotp_base_opp}) \\
\midrule
SOTP-High & $\Conf$ & IND-CONF & Groth16 ZK parfaite
  & $\Conf = 1$ \\
SOTP-High & $\Relver$ & BIND-REL & Groth16 (q-PKE)
  & $\geq 1-\negl(\lambda)$ \\
SOTP-High & $\Opp$ & -- & --
  & aucune (Rem.~\ref{rem:heuristique_sotp_high}) \\
\bottomrule
\end{tabular}
\end{table}
\cmtFormule{Tableau r\'ecapitulatif des neuf r\'eductions.}

\subsection{Corollaire central~: produit CLO game-based}

\begin{corollary}[Produit CLO depuis les jeux (section~\ref{sec:sotp_construction})]\label{cor:ann_cro_from_games}
Si IND-CONF, BIND-REL et EXT-OPP sont gagn\'es avec avantages respectifs au plus
$\varepsilon_{\mathsf{ind}}$, $\varepsilon_{\mathsf{bind}}$, $\varepsilon_{\mathsf{ext}}$, alors
\begin{equation}
\Conf_{\mathrm{calc}} \cdot \Relver \cdot \Opp
\geq \!\left(1-2\varepsilon_{\mathsf{ind}}\right)\!
\left(1-\varepsilon_{\mathsf{bind}}\right)\!
\left(1 - \frac{h(\varepsilon_{\mathsf{ext}})
  + \varepsilon_{\mathsf{ext}}\log_2(|\mathcal{V}_\lambda|-1)}
  {H(V|C)}\right)
\label{eq:ann_cro_product_games}
\end{equation}
\cmtFormule{Produit des trois bornes inf\'erieures.}
o\`u $\Conf_{\mathrm{calc}} = 1 - 2\varepsilon_{\mathsf{ind}}$ d\'esigne la confidentialit\'e
calculatoire. Pour $\varepsilon_{\mathsf{ind}},\varepsilon_{\mathsf{bind}},
\varepsilon_{\mathsf{ext}} \ll 1$, il vient
\begin{equation}
\Conf_{\mathrm{calc}} \cdot \Relver \cdot \Opp
\geq 1 - O\!\left(\varepsilon_{\mathsf{ind}}
  + \varepsilon_{\mathsf{bind}}
  + \varepsilon_{\mathsf{ext}}\log_2|\mathcal{V}_\lambda| + h(\varepsilon_{\mathsf{ext}})\right)
\label{eq:ann_cro_approx}
\end{equation}
\cmtFormule{Approximation au premier ordre pour petits $\varepsilon$.}
\end{corollary}

\begin{Preuve}
\textbf{Terme $\Conf_{\mathrm{calc}}$.} D\'efinition, avec le jeu IND-CONF
(Remarque~\ref{rem:ann_ind_conf_link}).
\textbf{Terme $\Relver$.} La Remarque~\ref{rem:ann_bind_rel_link} donne
$\Relver \geq 1-\varepsilon_{\mathsf{bind}} - \negl(\lambda)$.
\textbf{Terme $\Opp$.} Remarque~\ref{rem:ann_ext_opp_link}
(in\'egalit\'e de Fano, \eqref{eq:ann_fano_opp}).
Le produit $(1-a)(1-b)(1-c) \geq 1-a-b-c$ pour $a,b,c \in [0,1]$
donne~\eqref{eq:ann_cro_approx}.
\end{Preuve}

\begin{remark}[Pourquoi ce corollaire ne contredit pas le Trilemme]
\label{rem:ann_pas_de_contradiction}
Le Trilemme CRO (Th\'eor\`eme~\ref{thm:trilemme_cro_conditionnel}) majore le produit
$\Conf\cdot\Relrec\cdot\Opp$, o\`u $\Conf$ est statistique et $\Relrec$ mesure la
reconstruction du message. Le corollaire minore un produit d'une autre nature, qui fait intervenir la
confidentialit\'e \emph{calculatoire} et la fiabilit\'e de \emph{v\'erification}.
Les deux \'echelles ne se comparent pas. Un protocole peut avoir
$\Conf_{\mathrm{calc}}\cdot\Relver\cdot\Opp$ proche de~1 tout en ob\'eissant au
Trilemme, puisque $\Relrec$ reste alors de l'ordre de $2^{-\lambda}$
(Remarque~\ref{rem:rel_ver}). L'hypoth\`ese que
$\varepsilon_{\mathsf{ext}}$ est petit est pr\'ecis\'ement ce que le Trilemme rend
co\^uteux lorsque $\Conf$ statistique est proche de~1~; la
Remarque~\ref{rem:activite_delta} indique quelle quantit\'e ($\Delta$ ou Fano) limite
alors $\Opp$. La formulation initiale de ce paragraphe annon\c{c}ait un
\guillemotleft~encadrement~\guillemotright{} du produit par $\Delta/(\alpha\lambda)$ et une
constante $c \in [0{,}45;\,0{,}85]$~; la pr\'esente formulation distingue les deux \'echelles.
\cmt{Les deux produits rel\`event d'\'echelles distinctes et n'admettent pas d'encadrement commun.}
\end{remark}
\end{sloppypar}

% ============================================================================
\section{LIMITATIONS ET PORT\'EE DES R\'EDUCTIONS}
\addcontentsline{toc}{section}{Limitations et port\'ee}
\label{sec:gamebased_limitations}
% ============================================================================
\begin{sloppypar}

\begin{cmtEtape}{Objectif de cette section.}
\'Enum\'erer les cinq limitations des r\'eductions game-based, \`a savoir la
conditionnalit\'e, l'interpr\`ete honn\^ete, la r\'eduction one-way, le terme $1/q$
et la relation avec Lean4.
\end{cmtEtape}

\begin{enumerate}
\item \textbf{Conditionnalit\'e aux hypoth\`eses calculatoires.}
      La s\'ecurit\'e de SOTP-Low est prouv\'ee sous PRG et DLOG.
      SOTP-High via Groth16 repose sur la difficult\'e des
      couplages sur BN254, non post-quantique. Une extension aux
      primitives r\'esistantes aux quantiques (STARKs, r\'eseaux
      euclidiens) est une direction de recherche ouverte (section~\ref{sec:sotp_construction}).
      \cmt{La s\'ecurit\'e d\'epend d'hypoth\`eses non prouv\'ees (PRG, DLOG, couplages).}

\item \textbf{Mod\`ele d'interpr\`ete honn\^ete dans EXT-OPP.}
      Le jeu suppose que l'interpr\`ete tente sinc\`erement d'extraire
      le bon verdict. Un interpr\`ete \emph{adversarial} (juge
      corrompu) n'est pas mod\'elis\'e. Des extensions game-th\'eoriques
      (r\'egulateur vs entit\'e audit\'ee) sont identifi\'ees comme
      direction future (section~\ref{sec:sotp_construction}).
      \cmt{Un juge corrompu n'est pas couvert par le mod\`ele actuel.}

\item \textbf{R\'eduction SOTP-Base one-way.}
      Le Th\'eor\`eme~\ref{thm:ann_sotp_base_opp} garantit seulement l'implication
      \textit{dimension suffisante $\Rightarrow$ pr\'eservation JL}. Il ne garantit
      pas que la pr\'eservation des distances permette d'extraire le verdict,
      car cela d\'epend de la fonction $f_\Delta$ choisie par l'institution, dont ni
      l'existence ni la calculabilit\'e efficace ne sont \'etablies.
      \cmt{La pr\'eservation JL est n\'ecessaire, pas suffisante, pour l'opposabilit\'e.}

\item \textbf{Terme $1/q$ dans les bornes SOTP-Low.}
      Les \'equations~\eqref{eq:ann_sotp_low_conf}
      et~\eqref{eq:ann_sotp_low_rel} contiennent $1/q$
      (probabilit\'e de deviner l'exposant). Il est n\'egligeable en $\lambda$ au
      sens strict d\`es que l'ordre du groupe v\'erifie $\log_2 q \geq \lambda$, ce
      qui est le cas des param\`etres utilis\'es ($q \approx 2^{252}$, soit
      $1/q \approx 10^{-76}$). Le terme est conserv\'e pour que les bornes restent
      exactes \`a $\lambda$ fix\'e.
      \cmt{Le terme $1/q$ est n\'egligeable d\`es que $\log_2 q \geq \lambda$.}

\item \textbf{Relation avec la formalisation Lean4 (Annexe~\ref{chap:annexe_lean4}).}
      Les deux annexes se compl\`etent de la mani\`ere suivante.
      \begin{itemize}
      \item \textbf{Lean4 (Annexe~\ref{chap:annexe_lean4})} certifie m\'ecaniquement la
            d\'erivation alg\'ebrique de la borne
            $\Conf\cdot\Relrec\cdot\Opp \leq \Delta/\lambda + K/\lambda^2 + \varepsilon(\lambda)$
            (Th\'eor\`eme~\ref{thm:conditional_CRO_trilemma}), en posant Fano et
            Pinsker comme propri\'et\'es de chaque instance et l'extraction born\'ee
            comme hypoth\`ese
            \cmt{Certification m\'ecanique de la logique interne.}
      \item \textbf{Annexe~\ref{chap:annexe_gamebased} (ici)} \'etablit les
            propri\'et\'es cryptographiques de la famille SOTP (confidentialit\'e
            calculatoire, fiabilit\'e de v\'erification) sous PRG, DLOG et Groth16~;
            elle n'instancie pas les hypoth\`eses du Trilemme, qui portent sur
            $\Conf$ statistique et $\Relrec$
            \cmt{Instanciation dans des primitives calculatoires concr\`etes.}
      \end{itemize}
      Les deux niveaux sont donc compl\'ementaires sans se composer en une preuve
      unique, puisque le niveau \emph{formel} certifie la logique de la d\'erivation du
      Trilemme tandis que le niveau \emph{game-based} justifie les garanties
      cryptographiques des constructions.
      \cmt{La preuve se d\'eploie sur deux niveaux, formel (Lean4) et game-based (ici).}
\end{enumerate}
\end{sloppypar}